% !TeX root = ../../tesi.tex

\section{Il commercio agentico nella visione di Halue}
\label{sec:commercio-agentico}

Il \gloss{commercio-agentico}{commercio agentico} è l’insieme delle esperienze d’acquisto in cui un
agente di intelligenza artificiale accompagna il cliente nella ricerca, nella scelta e
nell’acquisto di un prodotto. Si sta sviluppando in due direzioni:

\begin{itemize}
	\item il negozio si apre ad assistenti generici di terze parti, che lo raggiungono con protocolli
	come MCP e UCP, e il cliente conversa altrove mentre il negozio fornisce catalogo e pagamento;
	\item l’assistente appartiene al marchio e opera all’interno del negozio, con la sua interfaccia
	e con strumenti scritti per il suo catalogo.
\end{itemize}

Le piattaforme promuovono la prima direzione, che un commerciante può adottare senza alcuno
sviluppo dedicato. La seconda richiede un progetto su misura, ed è quella in cui un’azienda come
Halue può offrire un servizio ai propri clienti. ChatCart la esplora.

Prima dello \emph{stage} Halue disponeva già di alcune premesse, cioè i negozi realizzati per
marchi che vendono su più mercati e gli agenti collegati ai dati dei negozi con strumenti propri
(\secref{sec:innovazione}). Non aveva invece esperienza su un assistente rivolto al cliente finale,
che pone problemi diversi da quelli di un agente a uso interno perché deve integrarsi con il tema
del negozio, rispondere in pochi secondi e comportarsi in modo prevedibile con qualsiasi
interlocutore. Lo \emph{stage} aveva lo scopo di costruire questa competenza su un caso
sufficientemente circoscritto da poter essere completato nella durata prevista.

Al termine dello \emph{stage}, il risultato costituisce per l’azienda una base di ricerca e
sviluppo. L’obiettivo dichiarato è arrivare a proporre questo tipo di soluzione ai clienti come
servizio, con tempi e forma ancora da stabilire. Il primo passo è stata la presentazione al
\emph{team} del 9 novembre 2026, l’ultimo giorno di \emph{stage}, articolata in trenta minuti
dedicati al prodotto, quindici alla parte tecnica e quindici alle domande. Le domande hanno
mostrato quali aspetti interessano a chi riprenderà il lavoro, in particolare le scelte di
progettazione, come la gestione dello stato nella sincronizzazione del carrello, e il modo in cui
ho istruito l’agente perché si comportasse come previsto.
